\section{Problem 650 --- ROUND 2}

\subsection*{1) ROUND-2 OBJECTIVE}
Path \textbf{(A)}: give a full proof of the target statement, and in fact of the stronger theorem
\[
 f(st)\le s+t
\qquad (s,t\ge 1).
\]
This implies the desired bound
\[
 f(m)\le 2\lceil \sqrt m\rceil,
\]
so the original square-root upper-bound question is answered affirmatively.

\subsection*{2) Round-1 FOUNDATION USED}
I use only one Round-1 item: the equivalent bipartite-matching reformulation.  Namely, for a set $A\subseteq \{1,\dots,N\}$ and an interval $I$ of length $2N$, one considers the bipartite graph with left side $A$, right side $I\cap\mathbb Z$, and edges $(a,b)$ whenever $a\mid b$.  Bounding the maximum matching is equivalent to bounding $f(m)$.

\subsection*{3) NEW INSIGHT / TOOL (ROUND-2)}
The new ingredient is an explicit Chinese-remainder construction which forces all integers in the interval divisible by members of $A$ to lie in only $s+t$ prescribed locations.  This is much stronger than the original target $f(m)\ll \sqrt m$.

\subsection*{4) ATTACK PLAN (ROUND-2)}
To prove $f(st)\le s+t$, it suffices to construct, for each $s,t\ge 1$,

\begin{itemize}
\item an integer $N$,
\item a set $A\subseteq \{1,\dots,N\}$ with $|A|=st$,
\item an interval $I$ of length $2N$,
\end{itemize}

such that at most $s+t$ integers of $I$ are divisible by at least one element of $A$.
Then any matching in the bipartite divisibility graph has size at most $s+t$.

The key steps are:
\begin{enumerate}
\item construct numbers $a_{i,j}$ whose pairwise gcd's are carefully controlled;
\item solve a compatible simultaneous congruence system so that each $a_{i,j}$ divides a prescribed integer $x_0+i$;
\item place the interval so that each $a_{i,j}$ has exactly two multiples inside it, and all such multiples lie in a fixed set of size $s+t$.
\end{enumerate}

\subsection*{5) WORK (ROUND-2)}
\paragraph{Theorem 650.1.}
For all integers $s,t\ge 1$,
\[
 f(st)\le s+t.
\]
Consequently, for every $m\ge 1$,
\[
 f(m)\le 2\lceil\sqrt m\rceil.
\]

\paragraph{Proof.}
We first prove the stronger statement $f(st)\le s+t$.
If $s=1$ or $t=1$, then any matching has size at most $|A|=st\le s+t$, so there is nothing to prove.
Hence assume from now on that
\[
 s,t\ge 2.
\]

Define
\[
 E:=\prod_{p\le st} p^{1+\lfloor \log_p(s-1)\rfloor},
 \qquad
 D:=E\cdot \mathrm{lcm}(1,2,\dots,st).
\]
Then for every prime $p\le st$ and every integer $\delta$ with
\[
1\le |\delta|\le s-1,
\]
one has
\[
 v_p(D)>v_p(\delta).
\]
Indeed, $v_p(\delta)\le \lfloor \log_p(s-1)\rfloor$, while the exponent of $p$ in $E$ is exactly $1+\lfloor \log_p(s-1)\rfloor$.

Now let
\[
\mathcal P:=\Bigl\{p>st:\ p\text{ prime and } p\mid (qD+\delta)
\text{ for some }1\le q\le t-1,
1\le |\delta|\le s-1\Bigr\}.
\]
This is a finite set of primes.
For each $p\in\mathcal P$, consider the residue classes
\[
 i-jD \pmod p
 \qquad (1\le i\le s,\ 0\le j\le t-1).
\]
There are at most $st$ such classes, and $p>st$, so they do not exhaust all residue classes modulo $p$.
Hence we may choose a residue class $r_p\pmod p$ such that
\[
 r_p\not\equiv i-jD \pmod p
 \qquad (1\le i\le s,\ 0\le j\le t-1).
\]
By the Chinese remainder theorem, there exists an integer $a$ such that
\[
 a\equiv r_p \pmod p
 \qquad (\forall p\in\mathcal P),
\]
and also
\[
 a>2(t-1)D+4s.
\]

For $1\le i\le s$ and $0\le j\le t-1$, define
\[
 a_{i,j}:=a+jD-i.
\]
These integers are pairwise distinct.  Indeed, if
\[
 a+jD-i=a+\ell D-k,
\]
then
\[
 (j-\ell)D=i-k.
\]
Since
\[
 |i-k|<s\le st\le D,
\]
we must have $j=\ell$, and then $i=k$.

Set
\[
 A:=\{a_{i,j}:1\le i\le s,\ 0\le j\le t-1\},
 \qquad
 N:=a+(t-1)D-1.
\]
Then $|A|=st$, and every element of $A$ lies in $\{1,\dots,N\}$ because
\[
 1\le a-s\le a_{i,j}\le a+(t-1)D-1=N.
\]

\paragraph{Claim 650.2.}
For any two pairs $(i,j)$ and $(k,\ell)$,
\[
 \gcd(a_{i,j},a_{k,\ell})\mid (k-i).
\]

\paragraph{Proof.}
Let
\[
 d:=\gcd(a_{i,j},a_{k,\ell}).
\]
Since
\[
 a_{i,j}-a_{k,\ell}=(j-\ell)D+(k-i),
\]
we have
\[
 d\mid (j-\ell)D+(k-i).
\]
There are three cases.

If $i=k$, then trivially $d\mid 0=(k-i)$.
If $i\ne k$ and $j=\ell$, then immediately $d\mid (k-i)$.
So only the case
\[
 i\ne k,
 \qquad
 j\ne \ell
\]
remains.
Let
\[
 q:=|j-\ell|\in\{1,\dots,t-1\}.
\]
Choose
\[
 \delta=
 \begin{cases}
 k-i,& j>\ell,\\
 i-k,& j<\ell.
 \end{cases}
\]
Then
\[
 1\le |\delta|=|k-i|\le s-1,
\]
and from the divisibility above we obtain
\[
 d\mid qD+\delta.
\]
We first show that every prime divisor of $d$ is at most $st$.
Suppose, for contradiction, that some prime $p>st$ divides $d$.
Because $p\mid qD+\delta$ with $1\le q\le t-1$ and $1\le |\delta|\le s-1$, the prime $p$ lies in $\mathcal P$ by definition.
On the other hand, $p\mid d\mid a_{i,j}=a+jD-i$, so
\[
 a\equiv i-jD \pmod p,
\]
contradicting the choice of $a\equiv r_p\pmod p$ with $r_p$ avoiding all residues $i-jD$.
Thus every prime divisor of $d$ is at most $st$.

Now let $p\le st$ be any prime divisor of $d$.
Since $d\mid qD+\delta$, we have $p\mid qD+\delta$.
Also,
\[
 v_p(qD)\ge v_p(D)>v_p(\delta).
\]
Therefore
\[
 v_p(qD+\delta)=v_p(\delta)=v_p(k-i).
\]
Since $d\mid qD+\delta$, it follows that
\[
 v_p(d)\le v_p(k-i).
\]
This holds for every prime divisor $p$ of $d$, hence indeed
\[
 d\mid (k-i).
\]

\hfill$\square$

By Claim 650.2, the congruence system
\[
 x\equiv -i \pmod{a_{i,j}}
 \qquad (1\le i\le s,\ 0\le j\le t-1)
\]
is pairwise compatible, because the difference between the residues $-i$ and $-k$ is $k-i$, and Claim 650.2 says that this is divisible by $\gcd(a_{i,j},a_{k,\ell})$.
By the generalized Chinese remainder theorem, there exists a solution $x_0$.
Replacing $x_0$ by a sufficiently large positive integer in the same residue class if necessary, we may assume that the interval introduced below lies inside $[1,\infty)$.

Choose an integer $T$ such that
\[
 2\bigl((t-1)D+s-1\bigr)<T\le a-2s.
\]
Such a choice is possible because
\[
 a>2(t-1)D+4s.
\]
Now define the interval
\[
 I:=(x_0-T,\ x_0-T+2N].
\]
It has length exactly $2N$.

Fix $(i,j)$ with $1\le i\le s$ and $0\le j\le t-1$.
From
\[
 x_0\equiv -i \pmod{a_{i,j}}
\]
we get
\[
 a_{i,j}\mid (x_0+i).
\]
Also, because $a_{i,j}=a+jD-i$,
\[
 x_0+a+jD=(x_0+i)+a_{i,j},
\]
so $x_0+a+jD$ is another multiple of $a_{i,j}$.
We claim that these are the \emph{only} two multiples of $a_{i,j}$ in $I$.

First, both lie in $I$.
The offset of $x_0+i$ from the left endpoint $x_0-T$ is $T+i$, which is positive and satisfies
\[
 T+i\le T+s\le a-s<2N.
\]
Hence $x_0+i\in I$.
Likewise, the offset of $x_0+a+jD$ from the left endpoint is $T+a+jD$, and
\[
 T+a+jD
 \le T+a+(t-1)D
 \le (a-2s)+a+(t-1)D
 <2N,
\]
because
\[
 2N=2a+2(t-1)D-2
\]
and
\[
 2N-\bigl((a-2s)+a+(t-1)D\bigr)=(t-1)D+2s-2>0.
\]
So $x_0+a+jD\in I$.

Next, the preceding multiple of $a_{i,j}$ is
\[
 (x_0+i)-a_{i,j}.
\]
Its offset from the left endpoint is
\[
 T+i-a_{i,j}.
\]
But
\[
 a_{i,j}=a+jD-i\ge a-s,
\]
while
\[
 T+i\le T+s\le a-s.
\]
Hence
\[
 T+i-a_{i,j}\le 0,
\]
so this preceding multiple is not in $I$.

Finally, the next multiple after $x_0+a+jD$ is
\[
 (x_0+i)+2a_{i,j}.
\]
Its offset from the left endpoint is
\[
 T+i+2a_{i,j}
 \ge T+1+2(a-s).
\]
Using
\[
 T>2\bigl((t-1)D+s-1\bigr)=2(t-1)D+2s-2,
\]
we obtain
\[
 T+i+2a_{i,j}
 >2(t-1)D+2s-2+1+2a-2s
 =2a+2(t-1)D-1
 >2N,
\]
since $2N=2a+2(t-1)D-2$.
Thus this next multiple also lies outside $I$.
Therefore $a_{i,j}$ has exactly two multiples in $I$, namely
\[
 x_0+i
 \qquad\text{and}\qquad
 x_0+a+jD.
\]

It follows that every integer in $I$ divisible by some element of $A$ lies in the set
\[
 B:=\{x_0+i:1\le i\le s\}\cup\{x_0+a+jD:0\le j\le t-1\}.
\]
The two parts of this union are disjoint, because an equality
\[
 x_0+i=x_0+a+jD
\]
would imply $i=a+jD$, impossible since $1\le i\le s<a\le a+jD$.
Hence
\[
 |B|=s+t.
\]
So the interval $I$ contains at most $s+t$ integers divisible by members of $A$.
Consequently the maximum matching size in the bipartite divisibility graph is at most $s+t$.
Therefore
\[
 f(st)\le s+t.
\]
This proves the stronger theorem.

To deduce the original one-parameter bound, let
\[
 n:=\lceil\sqrt m\rceil.
\]
Then $m\le n^2$.
Apply the stronger theorem with $s=t=n$ to obtain a set $A\subseteq\{1,\dots,N\}$ of size $n^2$ and an interval of length $2N$ for which at most $2n$ integers in the interval are divisible by elements of $A$.
Choose any subset $A'\subseteq A$ of size exactly $m$.
The same interval still contains at most $2n$ integers divisible by elements of $A'$, because $A'\subseteq A$.
Hence
\[
 f(m)\le 2n=2\lceil\sqrt m\rceil.
\]
This completes the proof.

\hfill$\square$

\subsection*{6) ADVERSARIAL VERIFICATION}
I checked the following potential failure points.

\begin{itemize}
\item \textbf{Pairwise distinctness of the moduli $a_{i,j}$.}  This is guaranteed by the inequality $|i-k|<D$.
\item \textbf{Existence of $a$.}  The exceptional-prime set $\mathcal P$ is finite, so CRT applies; adding a common multiple of all primes in $\mathcal P$ makes $a$ arbitrarily large.
\item \textbf{Existence of $T$.}  The interval
\[
\bigl(2((t-1)D+s-1),\ a-2s\bigr]
\]
is nonempty because $a>2(t-1)D+4s$.
\item \textbf{Generalized CRT step.}  Claim 650.2 gives exactly the compatibility condition required for a simultaneous solution of congruences with non-coprime moduli.
\item \textbf{Exactly two multiples in $I$.}  Both the previous multiple and the next multiple were checked explicitly by offsets from the left endpoint; the half-open interval notation avoids endpoint ambiguity.
\item \textbf{Passage from $n^2$ to $m$.}  Taking an $m$-element subset cannot create new divisible integers, so the one-parameter bound follows immediately.
\end{itemize}

No gap remains.

\subsection*{7) FINAL}
\textbf{FULL SOLUTION --- FULL PROOF.}

The original question is answered affirmatively.  In fact the stronger theorem
\[
 f(st)\le s+t
\qquad (s,t\ge 1)
\]
holds, and therefore
\[
 f(m)\le 2\lceil\sqrt m\rceil
\qquad (m\ge 1).
\]

\subsection*{8) COMPLETION ESTIMATE (MANDATORY)}
COMPLETION: 100\%

\subsection*{9) REFERENCES}
\begin{enumerate}
\item T. F. Bloom, \emph{Erd\H{o}s Problem \#650 and discussion thread}, accessed 2026-03-07.
\item Q. Tang, \emph{ep-650.tex} (GitHub proof source linked from the discussion thread), accessed 2026-03-07.
\end{enumerate}

\subsection*{10) OUTPUT FORMAT}
This section is part of the combined drop-in file \texttt{643\_644\_650\_solutions.tex}; no preamble is used.
